\part{Apéndices}
\appendix
\chapter{Problemas favortitos}
\begin{problem}
Encuentre todos los enteros positivos $n\geq3$ tales que existe
en el plano cartesiano un polígono convexo equilado de $n$ vértices todos
con coordenadas enteras.
\end{problem}

\begin{problem}
En la fiesta de San Juan del CTN,
cada grupo presenta un juego.
Uno de los juegos consiste en que
inicialmente se colocan $n \geq 1$ llaveritos en
una mesa y juegan entre sí
Mateo y Sofía, que son el que atiende en el
juego y un retador respectivamente.
Ambos juegan entre sí de forma alternada.
En su turno el jugador
debe quitar de $1$ a $k$ fichas
de la mesa.
Pierde el que quita la última ficha.

Para cada entero positivo $n$,
determine todos los enteros positivos
$k$ para los cuales Mateo
puede asegurar ganar el juego
sabiendo que éste empieza.
\begin{hints}
	\begin{hint}
		¿Qué pueden forzar alguno de los jugadores?
	\end{hint}
	\begin{hint}
		Múltiplos de $k + 1$.
	\end{hint}
	\begin{hint}
		Inducción en la cantidad de llaveritos.
	\end{hint}
\end{hints}
\solutionlabel{problem:ctnsanjuan-1}
\end{problem}
\begin{solutiontoproblem}{problem:ctnsanjuan-1}
	Probaremos que Mateo gana si y solo si $k + 1 \nmid n - 1$
	Procedemos por inducción en la cantidad de fichas.
	\paragraph{Caso 1 (caso base). $n \leq k + 1$. }
	Si $n = 1$, entonces Mateo pierde.
	Si $1 < n \leq k + 1$,
	entonces Mateo puede quitar $s \leq k$
	fichas para dejar exactamente una ficha,
	forzando a que Sofía quite la última ficha y pierda.

	En conclusión, para $n \leq k + 1$, Mateo pierde
	si y solo si $k + 1 \nmid n - 1$.
	Que es lo mismo que decir que
	Mateo pierde si y solo si
	\[n = t(k + 1) + 1\]
	con $t = 0$ y $t(k + 1) + 1 \leq n \leq (t + 1)(k + 1)$.

	\paragraph{Caso 2 (paso inductivo). $(t + 2)(k + 1) \geq n > (t + 1)(k + 1)$.}
	Si $n = (t + 1)(k + 1) + 1$,
	entonces si Mateo quita cualquier cantidad de llaveritos,
	entonces quedarán en la mesa $t(k + 1) + 1 < s \leq (t + 1)(k + 1)$
	fichas, y con todas esos valores sabemos Sofía puede ganar
	ya que puede quitar una cantidad de fichas para que queden
	en la mesa exactamente $t(k + 1) + 1$ fichas,
	por lo que Mateo pierde por hipótesis inductiva.

	Si $n > (t + 1)(k + 1) + 1$, entonces Mateo puede quitar $1 \leq s \leq k$
	llaveritos de la mesa para que queden exactamente $(t + 1)(k + 1) + 1$
	de ellos, cantidades en las cuales, probamos recién, Sofía
	pierde, o sea que Mateo gana, y la inducción queda completa.

	\begin{remark}
		En realidad, el jugador que gana, después de
		su primera jugada, solo tiene que quitar $(k + 1) - a$
		llaveritos el resto de sus jugadas,
		siendo $a$ la cantidad de llaveritos que
		jugó el otro jugador en el turno anterior.
	\end{remark}
\end{solutiontoproblem}

\begin{problem}
Se disponen dos tazas, una con leche y otra con té,
y ambas tienen la misma cantidad de leche y de té.
A continuación se pone una cucharadita de la taza de
leche a la taza de té. Después se pone una cucharadita
de la taza de té en la taza de leche.
¿Qué hay más, leche en la taza de té o té en la taza de leche?
\end{problem}
